\section{Relative elimination of source top handles}
\label{sec:relative}

\subsection{The source handle core}

The source-parallelity decomposition divides the metric cut into repeated
product pieces and a bounded exceptional core. Repeated normal disks, face
strips, and edge caps are recorded as affine interval families. Exceptional
tetrahedron chambers give zero-handles; exceptional paired-face regions give
one-handles; extreme gaps of interior source edges give two-handles; and
interior source vertices give three-handles. All attaching records retain
source tetrahedron, face, edge, and rank flags.

For the current constructor the verified inventory bounds are
\begin{equation}
 h_0\leq 6t,\qquad h_1\leq 8t,\qquad
 m=h_2\leq 12t,\qquad q=h_3\leq 4t,\qquad
 |\ports|\leq 66t.
 \label{eq:core-bounds}
\end{equation}
These bounds concern explicit exceptional objects; the represented
parallelity pieces can have much larger multiplicity and topology. Their
complete interfaces, not merely their homeomorphism types, are part of the
source packet~\cite{repo}.

Choose the canonical orientation of each global source edge $e$. Its two
oriented endpoint occurrences are $a_e,b_e$, even if they represent the same
global vertex. Let $w_e=|S_x\cap e|$. The $w_e$ intersection points divide
$e$ into gaps labelled $0,1,\ldots,w_e$. An interior source edge contributes
one exceptional two-handle if $w_e=0$ and two if $w_e>0$. Interior gaps
$1,\ldots,w_e-1$ belong to the selected product region.

For a source vertex $v$, put $\iota(v)=v$ when $v$ is interior and
$\iota(v)=\rho$ when it lies on the original boundary. The symbol $\rho$
also denotes an end on the new normal-cut boundary.

\begin{lemma}[Exact cap incidences]
\label{lem:caps}
For an interior source edge $e$, the two belt-disk end occurrences of its
exceptional two-handles are
\[
\begin{array}{c|c|c}
 \text{condition}&\text{gap rank}&\text{cap endpoints}\\\hline
 w_e=0&0&(\iota(a_e),\iota(b_e))\\
 w_e>0&0&(\iota(a_e),\rho)\\
 w_e>0&w_e&(\rho,\iota(b_e)).
\end{array}
\]
The two occurrences are retained with multiplicity. In particular, a
zero-weight loop at an interior vertex has two contacts with its
three-handle.
\end{lemma}

\begin{proof}
Before taking any quotient of source flags, take a small product
neighbourhood of each open edge gap. Transversely it is a disk; longitudinally
it is the indicated subinterval of the oriented edge. The transverse disk
times this interval is its dual two-handle. An endpoint at a source vertex
meets the vertex neighbourhood, which is a three-handle exactly when the
vertex link is a sphere. A boundary vertex has a disk link and contributes
an exterior end instead.

An endpoint at an intersection with a normal disk lies on a boundary sheet
of the metric cut. Thus a zero-weight edge keeps both source-vertex ends,
whereas each extreme gap of a positive-weight edge has one source-vertex end
and one cut-boundary end. The intervening gaps have two normal-sheet ends
and occur in the product families. The source edge-link and orientation
checks ensure that this description descends through all face identifications.
A global equality $a_e=b_e$ does not identify the two oriented end flags.
The argument is local to the split sheets, so a one-sided global normal
component causes no exception.
\end{proof}

\begin{figure}[htbp]
\centering
\begin{tikzpicture}[x=1cm,y=1cm,>=Stealth,
  vertex/.style={circle,draw=Ink,fill=white,inner sep=2.3pt},
  root/.style={circle,draw=Accent,fill=Accent!12,inner sep=2.3pt}]
 \node[vertex,label=below:$a$] (a) at (0,0) {};
 \node[vertex,label=below:$b$] (b) at (3,0) {};
 \draw[thick,Ink] (a)--node[above]{$h(e,0)$}(b);
 \node at (1.5,0.75) {$w_e=0$};
 \node[vertex,label=below:$a$] (c) at (6,0) {};
 \node[root,label=below:$\rho$] (r) at (8.2,0) {};
 \node[vertex,label=below:$b$] (d) at (10.4,0) {};
 \draw[thick,Ink] (c)--node[above]{$h(e,0)$}(r);
 \draw[thick,Ink] (r)--node[above]{$h(e,w_e)$}(d);
 \node at (8.2,0.75) {$w_e>0$};
\end{tikzpicture}
\caption{The cap graph records the extreme gaps. The common drawing of
 $\rho$ is a scheduling convention; distinct boundary pieces are never
 identified in the manifold. Boundary source vertices also map to $\rho$.}
\label{fig:caps}
\end{figure}

\subsection{The relative geometric hypothesis}

A graph trace alone cannot justify a marked three-manifold operation. The
entire protected interface must avoid the support of the cancellation.
Even the assertion that ports lie on zero- and one-handles is insufficient:
a three-handle attaching sphere can meet patches of those lower handles.
The native argument uses the actual rank charts.

\begin{lemma}[Separation in the source rank charts]
\label{lem:separation}
In the source-parallelity decomposition, the labelled interface rectangles
$\ports$, with compatible corner collars, can be realized disjointly from
the interiors and attaching regions of the exceptional two/three-handle
cancellation supports. This realization respects all source identifications
and the specified port charts.
\end{lemma}

\begin{proof}
Work first in the split tetrahedra, with their source flags kept distinct.
Choose mutually compatible small regular neighbourhoods of vertex sectors,
edge-gap sectors, and paired-face sectors. The normal disks bound the split
regions and separate their successive gap interiors. This is a finite PL
construction for each supplied $x$; it is used to prove the interpretation
of the compressed families, not to enumerate them.

A prism family has ranks $1\leq r<n$ between successive parallel disks of
one type. A face-strip family has ranks $1\leq r<a$ between successive
parallel normal arcs at its corner. An edge-cap family has ranks
$1\leq r<w_e$. These are strict interior ranks. The vertex sectors lie at
the source ends of the ordered gaps. The exceptional two-handles occupy
rank $0$, and also rank $w_e$ when $w_e>0$. Thus a selected product block
does not occupy a source vertex sector or an exceptional end-gap sector.
When an extreme gap and an interior product gap abut the same normal disk
before cutting, their closures lie on different local boundary copies in
the metric cut. This local separation also applies to a one-sided surface.

The core-facing free sides of a face-strip or edge-cap block are the
recorded port rectangles. Their source chamber and face flags are retained,
including for self-paired tetrahedra. If closures meet a lower-dimensional
corner during the initial subdivision, choose product collars on the two
flagged sides and round inside their respective sectors. Strict interior
rank places the port on the side separated from the end-gap/vertex
neighbourhood; no passage through a normal disk is required. Since the
paired rank maps agree along a source face, the collar choices can be made
equivariantly under each pairing. Edge and vertex links being a circle,
interval, sphere, or disk give precisely the local PL neighbourhoods needed
for simultaneous rounding.

After this compatible rounding the full port rectangles and their collars
miss the vertex balls and the selected exceptional two-handles, including
their lateral attaching annuli and their incident vertex-ball patches.
The proof uses the full source product blocks. It would not follow from a
list of port owner dimensions or a matrix of signed exponents.
\end{proof}

\begin{remark}
The protected set in \cref{lem:separation} is the native parallelity
interface. A future inherited boundary-pattern curve can cross a vertex
patch or an extreme-gap support even when it is carried by a zero/one
skeleton. Such a curve requires a separate avoidance or transport proof.
This article does not silently enlarge $\ports$ to an arbitrary pattern.
\end{remark}

\subsection{One cancellation and a complete rooted schedule}

\begin{lemma}[Relative single-contact cancellation]
\label{lem:relative-move}
Let $h$ be an exceptional two-handle with exactly one contact, counted with
multiplicity, with the currently remaining three-handles. Let that contact
be with $D$. Assume that the other cap of $h$ is on the current exterior
boundary and that the marked interface misses the cancellation support.
Removing $h$ and $D$ gives a homeomorphic manifold relative to that
interface. Each surviving two-handle keeps its original attaching cycle to
the zero/one-handle union. Its cap contacts with $D$ become exterior contacts.
\end{lemma}

\begin{proof}
This is a two/three-handle cancellation with geometric incidence one. In
the source handle charts, $h$ and $D$ are balls glued along their unique
common cap disk. Their union is a ball. Its attachment to the remaining
handles is a disk: the complementary disk in the attaching sphere of $D$
is extended by the lateral attaching annulus of $h$. The other cap is
exterior. Removing a ball attached along a boundary disk preserves the
homeomorphism type by adjusting a collar of that disk.

The collar can be chosen away from the protected interface. This gives a
homeomorphism fixed on the interface and its specified charts. It does not
assert that every point of every undeleted lower handle is fixed: the
attaching-disk collar must move. The reduced presentation itself consists
of the undeleted handles, whose mutual attaching records remain unchanged.
In particular, the other two-handle attaching cycles are literally the
same records. Only their upper cap status changes when $D$ is erased.
\end{proof}

Define the cap multigraph $G_x$ on the $q$ interior vertex three-handles and
the formal root $\rho$. Each exceptional two-handle is an edge with the
endpoint occurrences in \cref{lem:caps}. An exterior--exterior handle is a
loop at $\rho$. A loop at an interior vertex is also retained as a loop,
with two incidences.

\begin{lemma}[Rooted deletion criterion]
\label{lem:rooted}
For any such cap graph, all three-handles can be deleted by single-contact
steps if and only if every graph vertex reaches $\rho$. The edge sets of
complete schedules are exactly the spanning trees of $G_x$.
\end{lemma}

\begin{proof}
Given a spanning tree rooted at $\rho$, delete each vertex after its parent,
using the edge to that parent. At that instant the edge has exactly one live
endpoint. Conversely, orient each chosen edge from the vertex it deletes
toward its other end. That end is $\rho$ or a previously deleted vertex.
Following the chosen edges therefore reaches $\rho$ and never forms a
cycle. Exactly $q$ edges are selected on $q+1$ vertices, so they form a
spanning tree. A component not meeting $\rho$ cannot supply its first legal
step. A loop never has exactly one live endpoint.
\end{proof}

\begin{theorem}[Source-relative top-handle elimination]
\label{thm:top}
For every admissible normal vector in the stated source setting, all $q$
source three-handles can be removed by exactly $q$ exceptional two/three
cancellations. The resulting core $K'$ has no three-handles and satisfies
\[
 (K,\ports)\cong(K',\ports).
\]
All labelled port rectangles and charts, all zero/one-handle IDs, and all
surviving two-handle attaching cycles can be retained. Reattaching the same
parallelity pieces along the same interfaces gives a homeomorphic whole
cut.
\end{theorem}

\begin{proof}
Consider a component $U$ of the graph of zero-weight source edges on the
interior vertices. A positive-weight edge incident to $U$ gives an exterior
contact by \cref{lem:caps}. A zero-weight edge from $U$ to a boundary vertex
does the same. If neither occurred, every source edge incident to $U$ would
have weight zero and both ends in $U$. This would isolate $U$ in the source
one-skeleton, contrary to connectedness of that skeleton and nonemptiness
of the source boundary. Hence every such component reaches an exterior
contact, and $G_x$ is rooted-connected.

Use \cref{lem:rooted} for the schedule, \cref{lem:separation} for its protected
interface, and \cref{lem:relative-move} for each step. Induction gives the
relative homeomorphism. Boundary components represented by the same formal
root are never glued together: every actual cancellation has only one live
three-handle end. Finally, a homeomorphism fixed on the entire gluing
interface extends by the identity over every reattached parallelity piece.
\end{proof}

Deleting equal numbers of two- and three-handles preserves Euler
characteristic. Surviving two-handle IDs may be densely renumbered for the
old terminal consumer after source replay, but a certificate must record
the original IDs used by the trace. This distinction prevents an accidental
reassignment of attaching cycles when rows are removed.
